% Monografia proiectului, ediție de probă. Se construiește cu build.sh, care
% pregătește planșele, datele și versiune.tex în build/; aici doar textul.
\documentclass[11pt,a4paper,oneside]{article}
\input{versiune}   % \RAD (rădăcina depozitului), \versiune (commitul)

% ---------------------------------------------------------------- limba
\usepackage{babel}
\babelprovide[import,main]{romanian}
% Tiparele de despărțire sunt aceleași pe care le folosește convertorul.
\directlua{
  local f = assert(io.open("\luaescapestring{\RAD}/tools/converter/patterns/hyph-ro.pat.txt"))
  local s = f:read("*a"); f:close()
  lang.patterns(lang.new(tex.language), s)
}
\lefthyphenmin=2 \righthyphenmin=2

% ---------------------------------------------------------------- literele
\usepackage{fontspec}
\defaultfontfeatures{Ligatures=TeX}
\edef\libron{\noexpand\setmainfont{Libron_R}[Path=\RAD/tools/fonttool/fonts/libron/,
  Extension=.ttf, UprightFont=*-Regular, BoldFont=*-Bold,
  ItalicFont=*-Italic, BoldItalicFont=*-BoldItalic,
  Numbers={OldStyle,Proportional}]}
\libron
\newcommand\ohm{{\fontfamily{lmr}\selectfont Ω}}
\newcommand\lin{\addfontfeatures{Numbers={Lining,Monospaced}}}
\usepackage{amsmath}
\usepackage{unicode-math}
\setmathfont{Latin Modern Math}
% cifrele și literele latine din formule, tot în Libron; restul din Latin Modern
\edef\mathlibron{\noexpand\setmathfont{Libron_R-Regular.ttf}[Path=\RAD/tools/fonttool/fonts/libron/,
  range={up/{num,latin,Latin}}]%
  \noexpand\setmathfont{Libron_R-Italic.ttf}[Path=\RAD/tools/fonttool/fonts/libron/,
  range={it/{latin,Latin}}]}
\mathlibron
\usepackage[protrusion=true,expansion=true]{microtype}

% Paragrafele acestui text se rup cu parametrii dispozitivului (§ II).
\linepenalty=10 \hyphenpenalty=100 \exhyphenpenalty=100
\doublehyphendemerits=100000 \finalhyphendemerits=5000 \adjdemerits=10000
\tolerance=400 \emergencystretch=1em

% ---------------------------------------------------------------- pagina
\usepackage[a4paper, left=27mm, textwidth=118mm, marginparsep=7mm,
  marginparwidth=40mm, top=30mm, bottom=32mm, headsep=9mm, footskip=13mm]{geometry}
\newlength\latplin \setlength\latplin{\dimexpr\textwidth+\marginparsep+\marginparwidth}
\usepackage{fancyhdr}
\pagestyle{fancy}\fancyhf{}
\renewcommand\headrulewidth{0pt}
\fancyhead[C]{\footnotesize\scshape despre așezarea cărților pe hârtie electronică}
\fancyfoot[C]{\small\thepage}
\setlength\parindent{1.4em}

\usepackage{titlesec, titletoc, needspace}
\titlecontents{section}[2.2em]{\addvspace{0.3ex}}{\contentslabel[\scshape\thecontentslabel]{2.2em}}{}{\titlerule*[0.7pc]{.}\contentspage}
\newcommand\ornament{%
  \tikz[baseline=-0.6ex]{\draw[line width=0.3pt] (-1.1,0)--(-0.12,0) (0.12,0)--(1.1,0);
    \draw[line width=0.3pt] (0,-0.07)--(0.07,0)--(0,0.07)--(-0.07,0)--cycle;}}
\titleformat{\section}[display]{\needspace{14\baselineskip}\centering\normalfont}
  {\scshape §\,\Roman{section}}{2pt}{\large\itshape}[\vspace{3pt}\ornament]
\titlespacing*{\section}{0pt}{5ex}{3ex}
\renewcommand\thesection{\Roman{section}}
\renewcommand\thesubsection{\arabic{subsection}}
\titleformat{\subsection}[runin]{\normalfont\itshape}{\thesubsection.}{0.4em}{}[.]
\titlespacing*{\subsection}{\parindent}{1.2ex}{0.6em}

\usepackage{lettrine}
\renewcommand\LettrineFontHook{\normalfont}
\setcounter{DefaultLines}{3}

\usepackage{marginnote}
\renewcommand*{\raggedleftmarginnote}{\raggedright}
\newcommand\sursa[1]{\marginnote{\footnotesize\itshape\raggedright\hyphenpenalty=10000\exhyphenpenalty=10000 #1}}
\usepackage[hidelinks]{hyperref}
\urlstyle{same}
\newcommand\fisier[1]{{\lin\upshape\nolinkurl{#1}}}
\hypersetup{pdftitle={Despre așezarea cărților pe hârtie electronică},
  pdfauthor={Beleniuc Ilie}, pdflang={ro}}

% ---------------------------------------------------------------- figuri, tabele
\usepackage{graphicx, booktabs, tabularx, array, enumitem}
\usepackage{tikz}
\usetikzlibrary{calc, arrows.meta}
\usepackage{pgfplots, pgfplotstable}
\usepgfplotslibrary{groupplots}
\pgfplotsset{compat=1.18,
  every axis/.append style={axis lines*=left, tick align=outside,
    axis line style={line width=0.4pt}, tick style={line width=0.4pt, black},
    every tick label/.append style={font=\footnotesize\lin},
    label style={font=\small\itshape},
    legend style={draw=none, fill=none, font=\footnotesize},
    /pgf/number format/use comma, /pgf/number format/1000 sep={.}}}
\usepackage[labelsep=period, font=small, labelfont=sc, justification=justified,
  singlelinecheck=true]{caption}
\DeclareCaptionType[fileext=lop, placement=tp]{plansa}[Planșa][Lista planșelor]
\renewcommand\theplansa{\Roman{plansa}}
\addto\captionsromanian{\renewcommand\figurename{Fig.}\renewcommand\tablename{Tabelul}}
\renewcommand\figurename{Fig.}\renewcommand\tablename{Tabelul}
\setlist{itemsep=0.2ex, topsep=0.6ex}

% literă pe planșă, cu firul ei până la piesă
\newcommand\lit[3]{% eticheta, unde stă litera, unde arată
  \draw[line width=0.25pt] (#2) -- (#3);
  \fill (#3) circle (0.18);
  \node[fill=white, inner sep=0.6pt, font=\small\itshape] at (#2) {#1};}

\newcommand\ceva[1]{\IfFileExists{#1}}
\newcommand\lipseste[2]{\fbox{\parbox[c][#2][c]{#1}{\centering\footnotesize\itshape
  lipsește la această compilare\\(vezi README)}}}

\begin{document}
% ================================================================ titlul
\begin{titlepage}
\centering
\vspace*{8mm}
{\scshape universitatea „ștefan cel mare” din suceava}\\[1mm]
{\small\itshape Facultatea de Inginerie Electrică și Știința Calculatoarelor}

\vspace{26mm}
{\fontsize{27}{31}\selectfont Despre așezarea cărților\\[1mm] pe hârtie electronică}

\vspace{7mm}\ornament\vspace{7mm}

\begin{minipage}{0.78\textwidth}\centering\itshape
Cercetare asupra unui cititor de buzunar cu ecran electroforetic, asupra felului
în care își rupe paragrafele în rânduri, a încăperii cărților în memoria lui,
precum și a celor ce nu s-au putut încă măsura
\end{minipage}

\vspace{24mm}
{\small de}\\[2mm]
{\Large\scshape beleniuc ilie}

\vspace*{\fill}
{\small Proiect de licență\\[1mm] ediție de probă, \lin\versiune\\[1mm] noiembrie 2026}
\end{titlepage}

% ================================================================ rezumat
\thispagestyle{empty}
\begin{center}
\begin{minipage}{0.84\textwidth}
\begin{center}\scshape rezumat\end{center}
\small
Lucrarea descrie un cititor de cărți construit în jurul unui ecran electroforetic
de 3,7 țoli (416 × 240 de pixeli) și al unui microcontroler ESP32-S3, împreună cu
uneltele de pe calculator care îi pregătesc cărțile. Partea măsurată este
așezarea textului: pe o coloană de 216 pixeli, algoritmul lui Knuth și Plass, cu
despărțire în silabe făcută dinainte pe calculator, micșorează spațiul mediu
adăugat la justificare cu 27~\% și percentila 95 cu 34~\% față de umplerea lacomă
a rândurilor, la 127.876 de rânduri de proză. Un model al algoritmului din TeX,
verificat paragraf cu paragraf față de LuaTeX, arată că motorul dispozitivului dă
rupturile lui TeX la 98,63~\% din paragrafe, iar diferențele rămase sunt alese.
Formatul cărților ține textul comprimat la 27–46~\% din mărimea lui. Placa,
carcasa și măsurătorile de curent sunt proiectate, dar nimic nu a rulat încă pe
hardware; lucrarea spune, la fiecare cifră, de unde vine și cât e de sigură.

\medskip
{\raggedright\textit{Cuvinte-cheie:} hârtie electronică, ruperea paragrafelor,
Knuth–Plass, despărțire în silabe, compresie DEFLATE, ESP32-S3.\par}
\end{minipage}
\end{center}

\vspace{10mm}
\begin{center}
\begin{minipage}{0.84\textwidth}
\renewcommand\contentsname{\centering\normalsize\scshape cuprins}
\small
\setcounter{tocdepth}{1}
\tableofcontents
\end{minipage}
\end{center}
\clearpage

% ================================================================ I
\section{Introducere}

\lettrine{U}{n cititor de cărți} electronice pare un lucru rezolvat. Cele din
comerț arată textul bine, dar o fac cu un procesor de aplicații, un sistem de
operare întreg și sute de megaocteți de memorie. Întrebarea acestei lucrări e
mai îngustă: cât de bine se poate așeza o carte pe un ecran mic, cu un
microcontroler cu 512 KB de memorie internă, 8 MB de PSRAM și un ceas de 240 MHz, dacă
munca grea se face o singură dată, pe calculator.

Proiectul are patru părți. \emph{Convertorul} citește EPUB sau text simplu și
scrie un fișier IBV: textul despărțit în silabe, tăiat în blocuri comprimate,
cu un index care permite saltul la orice poziție. \emph{Uneltele de fonturi}
pregătesc literele ca imagini gata de pus pe ecran, la cinci mărimi.
\emph{Programul dispozitivului} așază paginile, le desenează și le trimite
panoului. \emph{Simulatorul} compilează pe calculator exact sursele de așezare ale
dispozitivului, așa că paginile din Planșa~II sunt, pixel cu pixel, cele pe care
le-ar arăta cititorul.
\sursa{sim/CMakeLists.txt compilează firmware/components/core}

\subsection{Ecranul}
Panoul este unul electroforetic. Între doi electrozi stă un strat de capsule
microscopice, fiecare cu particule albe și negre încărcate cu sarcini opuse
\cite{comiskey}. Tensiunea dintre electrodul comun și cel al unui pixel trage
particulele de un fel sub sticlă, pe celelalte la fund. Imaginea rămâne după ce
tensiunea dispare: pagina nu costă nimic cât timp e citită, ci doar când e
întoarsă. De aici vine măsura care contează pentru un asemenea cititor, energia
pe pagină întoarsă, nu puterea medie (§\,V).

\begin{plansa}[t]
\makebox[\textwidth][l]{%
\ceva{capsule.mps}{%
\begin{tikzpicture}
  \node[anchor=south west, inner sep=0] at (0,0) {\includegraphics[width=\latplin]{capsule.mps}};
  \begin{scope}[x=\latplin/143.2, y=\latplin/143.2, shift={(71.6,15.8)}]
    \foreach \cx/\st in {-49.5/alb, -4.9/negru, 44.6/alb}
      \node[font=\footnotesize\scshape] at (\cx, 17.6) {\st};
    \lit{a}{-66,18}{-66,10.6}
    \lit{g}{-55,-19.5}{-49.2,-4.5}
    \lit{b}{-44,-19.5}{-43.8,-6.6}
    \lit{c}{-22,-19.5}{-19.6,-5}
    \lit{e}{0,-19.5}{2,-9.3}
    \lit{d}{36,-19.5}{39.2,-5}
    \lit{f}{56,-19.5}{56,-11.5}
  \end{scope}
\end{tikzpicture}}{\lipseste{\latplin}{30mm}}}
\caption[Principiul panoului electroforetic]{Principiul panoului
electroforetic, în secțiune; nu la scară. \textit{a}, folia de deasupra, cu
electrodul comun transparent; \textit{b}, membrana unei capsule; \textit{c},
particule albe; \textit{d}, particule negre; \textit{e}, electrozii celor trei
pixeli; \textit{f}, substratul; \textit{g}, liantul dintre capsule. Pixelii arată,
de la stânga, alb, negru, alb. Capsula a cincea stă pe granița dintre doi pixeli
și se împarte: de aici marginile ușor neclare ale literelor pe asemenea ecrane.
Desenul arată principiul descris de Comiskey și alții \cite{comiskey}; stratul
panoului folosit nu a fost secționat.}
\end{plansa}

\subsection{Ce se arată aici}
Paragrafele de mai jos se opresc la ce s-a putut verifica pe calculator.
Fiecare cifră are pe margine fișierul din care vine; figurile sunt desenate la
compilare, din aceleași fișiere.

\begin{center}
\begin{minipage}{0.9\textwidth}
\rule{\linewidth}{0.4pt}\\[0.5ex]
{\small\scshape declarație asupra limitelor.}\enspace{\small
Nimic din această lucrare nu a rulat pe hardware. Placa proprie este proiectată
și verificată electric, dar nu a fost comandată; carcasa nu a fost tipărită;
modulul de măsurare a curentului nu a fost cumpărat. Programul dispozitivului
trece testele într-un emulator al lui ESP32-S3, nu pe cip. Cifrele de consum din
§\,V sunt valori tipice din foile de catalog, nu măsurători.}\\[-0.6ex]
\rule{\linewidth}{0.4pt}
\end{minipage}
\end{center}

% ================================================================ II
\section{Despre ruperea paragrafelor}

\lettrine{P}{e un rând îngust}, fiecare spațiu contează. Coloana de text a
cititorului are 216 pixeli; la mărimea implicită a literei încap pe ea cam
patru-cinci cuvinte. Când textul e aliniat la ambele margini, ce lipsește
rândului ca să ajungă la margine se împarte între spațiile lui, iar un rând cu
două spații și un gol mare de umplut se vede de departe.

\subsection{Măsura}
Fie $s_0$ lățimea spațiului simplu al fontului și $s$ lățimea spațiului după
justificare. Atunci
\begin{equation}
  r = \frac{s - s_0}{s_0}
\end{equation}
spune cât s-a lățit spațiul: $r = 0$ pentru un spațiu firesc, $r = 1$ pentru
unul dublu. Un spațiu se poate și îngusta, cu cel mult o treime, deci
$r \ge -\tfrac13$.
\sursa{firmware/components/ core/layout.c: STRETCH, SHRINK}

\subsection{Doi algoritmi}
\emph{Umplerea lacomă} pune pe fiecare rând câte cuvinte încap și trece mai
departe; e ce fac cele mai multe cititoare simple. \emph{Algoritmul lui Knuth și
Plass} \cite{knuthplass} privește tot paragraful deodată și alege, dintre toate
felurile de a-l rupe, pe cel cu penalizarea totală cea mai mică. Fiecare spațiu
are o elasticitate: $\sigma = 2s_0$ când se lățește și $\sigma = s_0/3$ când se
îngustează. Urâțenia unui rând este
\begin{equation}
  b = 100\,\lvert \rho\rvert^{3}, \qquad \rho = \frac{s - s_0}{\sigma},
\end{equation}
iar penalizarea rândului $j$ este
\begin{equation}
  d_j = (\lambda + b_j)^2 + \pi_j^2 + \alpha\,[\text{două cratime la rând}]
        + \gamma\,[\text{salt de clasă}],
\end{equation}
cu $\lambda = 10$, $\pi_j = 100$ dacă rândul se termină într-o cratimă,
$\alpha = 100\,000$ și $\gamma = 10\,000$ dacă rândul e cu cel puțin două clase
mai strâns sau mai larg decât cel de deasupra. Penultimul rând despărțit mai
costă 5\,000. Algoritmul încearcă întâi cu $b \le 400$, adică
$\rho \le 1{,}59$ sau $r \le 3{,}17$; dacă paragraful nu se poate rupe așa, acceptă orice rând care încape.
Parametrii sunt cei din TeX \cite{texprogram}, cu două schimbări făcute cu
intenție (§\,III); textul pe care îl citiți este rupt cu aceiași parametri.
\sursa{layout.c: LINE, HYPHEN, DOUBLE\_HYPHEN, FINAL\_HYPHEN, ADJ\_FITNESS, TOLERANCE}

Despărțirea în silabe nu se face pe dispozitiv. Convertorul aplică tiparele lui
Liang \cite{liang} pentru română și pune în text cratime moi (U+00AD), iar
dispozitivul doar le folosește. Costul în spațiu e măsurat în §\,IV.

\subsection{Experimentul}
Opt cărți în proză (Caragiale, Slavici, Regina Maria, Pușkin în rusă, Austen,
Carroll, Melville, Shelley) au fost paginate în patru condiții: cei doi
algoritmi, cu și fără despărțire, la fiecare din cele cinci trepte de mărime a
literei. Rezultatul e determinist: cinci rulări repetate dau aceleași pagini.
Înainte de rulare au fost scrise cinci predicții; patru au trecut.
\sursa{measurements/E6/ METODA.md, REZULTATE.md}

\begin{table}[t]
\centering\small
\caption{Proza, la treapta implicită (2): toate rândurile justificate ale celor
opt cărți, puse împreună. „Perechi” sunt rânduri consecutive încheiate amândouă
cu cratimă, la mia de rânduri.}
\label{tab:e6}
\lin\setlength\tabcolsep{4pt}
\begin{tabular}{@{}l rrrrr@{}}
\toprule
& \multicolumn{2}{c}{$r$} & rânduri & perechi & \\
\cmidrule(lr){2-3}
condiția & medie & p95 & $r > 1$ & ‰ & pagini \\
\midrule
lacom, fără despărțire        & 2,89 & 9,21 & 69,2 \% & 0    & 10.086 \\
Knuth–Plass, fără despărțire  & 2,13 & 5,68 & 67,6 \% & 0    & 9.884 \\
lacom, cu despărțire          & 1,53 & 4,23 & 55,3 \% & 47,6 & 9.755 \\
Knuth–Plass, cu despărțire    & 1,12 & 2,80 & 49,7 \% & 22,1 & 9.579 \\
\bottomrule
\end{tabular}
\sursa{measurements/E6/ tabele.md, tabelul 1}
\end{table}

\subsection{Ce arată}
Despărțirea în silabe contează cel mai mult: fără ea, spațiul mediu este de
aproape două ori mai mare, la ambii algoritmi, iar Knuth–Plass fără despărțire
(2,13) e mai rău decât umplerea lacomă cu despărțire (1,53). Peste despărțire,
Knuth–Plass scade spațiul mediu cu 27~\% și percentila 95 cu 34~\%, are de două
ori mai puține perechi de cratime și ocupă cu 1,8~\% mai puține pagini
(tabelul~\ref{tab:e6}). Coada distribuției se vede în fig.~\ref{fig:hist}: la
umplerea lacomă, 6,0~\% din rânduri au spații de aproape cinci ori mai late
decât cele firești ($r \ge 3{,}95$); la Knuth–Plass, 0,9~\%.

\begin{figure}[t]
\centering
\ceva{date/histograma.dat}{%
\begin{tikzpicture}
\begin{axis}[width=\textwidth, height=52mm, xmin=-0.5, xmax=3.95, ymin=0,
  xlabel={$r$}, ylabel={\% din rânduri}, clip=false,
  xtick={-0.5,0,1,2,3}, extra x ticks={0.333}, extra x tick labels={},
  legend pos=north east, legend cell align=left]
  \addplot[const plot mark left, draw=black!45, fill=black!10, line width=0.3pt]
    table[x=r, y=greedy, restrict x to domain=-1:3.92] {date/histograma.dat} \closedcycle;
  \addplot[const plot mark left, black, line width=0.7pt]
    table[x=r, y=kp, restrict x to domain=-1:3.92] {date/histograma.dat};
  \draw[densely dotted, line width=0.3pt] (axis cs:1,0) -- ({axis cs:1,0}|-{rel axis cs:0,0.92})
    node[above, font=\footnotesize\itshape] {spațiu dublu};
  \legend{umplerea lacomă, Knuth–Plass}
\end{axis}
\end{tikzpicture}}{\lipseste{\textwidth}{50mm}}
\caption{Distribuția lui $r$ pe rândurile justificate, cu despărțire, la
treapta 2, pe intervale de 0,05. Valorile negative sunt rândurile strânse, pe
care doar Knuth–Plass le folosește. Ultimul interval ($r \ge 3{,}95$) nu e
desenat.}
\label{fig:hist}
\sursa{E6/date/ histograma.csv}
\end{figure}

Litera mai mare strică spațierea la toate variantele (fig.~\ref{fig:trepte}):
de la treapta 0 la 4, spațiul mediu crește de 1,5–1,7 ori. Knuth–Plass cu
despărțire la treapta cea mai mare (1,45) e cam cât umplerea lacomă cu
despărțire la treapta 1 (1,39).

\begin{figure}[t]
\centering
\ceva{date/trepte.dat}{%
\begin{tikzpicture}
\begin{groupplot}[group style={group size=2 by 1, horizontal sep=16mm},
  width=0.5\textwidth, height=50mm, xmin=-0.3, xmax=4.3, xtick={0,...,4},
  xlabel={treapta}, ymin=0]
\nextgroupplot[ylabel={$r$ mediu}, ymax=4,
  legend to name=legtrepte, legend columns=2, legend cell align=left,
  legend style={/tikz/every even column/.append style={column sep=6mm}}]
  \addplot[densely dashed, mark=o, mark size=1.6pt] table[x=treapta, y=greedy_nu_m] {date/trepte.dat};
  \addplot[densely dashed, mark=*, mark size=1.6pt] table[x=treapta, y=kp_nu_m] {date/trepte.dat};
  \addplot[solid, mark=o, mark size=1.6pt] table[x=treapta, y=greedy_da_m] {date/trepte.dat};
  \addplot[solid, mark=*, mark size=1.6pt] table[x=treapta, y=kp_da_m] {date/trepte.dat};
  \legend{lacom fără despărțire, Knuth–Plass fără despărțire,
          lacom cu despărțire, Knuth–Plass cu despărțire}
\nextgroupplot[ylabel={$r$, percentila 95}, ymax=14]
  \addplot[densely dashed, mark=o, mark size=1.6pt] table[x=treapta, y=greedy_nu_p95] {date/trepte.dat};
  \addplot[densely dashed, mark=*, mark size=1.6pt] table[x=treapta, y=kp_nu_p95] {date/trepte.dat};
  \addplot[solid, mark=o, mark size=1.6pt] table[x=treapta, y=greedy_da_p95] {date/trepte.dat};
  \addplot[solid, mark=*, mark size=1.6pt] table[x=treapta, y=kp_da_p95] {date/trepte.dat};
\end{groupplot}
\end{tikzpicture}\\[1ex]
\pgfplotslegendfromname{legtrepte}}{\lipseste{\textwidth}{50mm}}
\caption{Spațierea pe treptele de mărime ale literei (0 e cea mai mică).}
\label{fig:trepte}
\sursa{E6/date/linii.csv}
\end{figure}

Predicția care nu a trecut privește timpul. Pe calculator, Knuth–Plass e de
2,32 ori mai lent decât umplerea lacomă, nu de cel mult două ori. Pe pagină,
asta înseamnă paginarea unei cărți întregi în 376 de milisecunde în loc de 162,
pe un laptop. Cât costă pe ESP32-S3 nu e măsurat.
\sursa{E6/date/timp.csv}

\begin{plansa}[p]
\centering
\newcommand\pag[2]{%
  \ceva{#1}{\begin{tikzpicture}
    \node[inner sep=0, anchor=south west] (p) at (0,0)
      {\includegraphics[width=48mm, interpolate=false]{#1}};
    \draw[line width=0.4pt] ($(p.south west)+(-2mm,-2mm)$) rectangle ($(p.north east)+(2mm,2mm)$);
    \draw[line width=1.2pt] ($(p.south west)+(-3.2mm,-3.2mm)$) rectangle ($(p.north east)+(3.2mm,3.2mm)$);
    #2
  \end{tikzpicture}}{\lipseste{48mm}{83.2mm}}}
\makebox[\textwidth][c]{%
\pag{pagini/lacom-4.png}{%
  \draw[overlay, line width=0.4pt, -{Stealth[length=2mm]}] (-9mm, 0.355*83.2mm) -- (-4.2mm, 0.355*83.2mm);}%
\hspace{9mm}%
\pag{pagini/kp-4.png}{}}

\vspace{4mm}
\begin{minipage}{48mm}\centering\small\itshape a. umplerea lacomă\end{minipage}\hspace{15.4mm}%
\begin{minipage}{48mm}\centering\small\itshape b. Knuth–Plass\end{minipage}
\caption[Aceeași pagină, ruptă în două feluri]{Pagina a patra din Caragiale,
\textit{Nuvele}, la mărime naturală pe ecranul de 3,7 țoli, așa cum o desenează
simulatorul cu codul dispozitivului. Ambele variante despart în silabe. Rândul
arătat de săgeată, „tut stăpânì\ldots{} a biruit-o”, nu are unde să ia o silabă
de pe rândul următor fără să iasă din margine, așa că umplerea lacomă îl
întinde; Knuth–Plass mută „pu-” un rând mai jos și împarte golul pe trei
rânduri. Pe toată cartea, varianta \textit{b} are 526 de pagini în loc de 537.}
\end{plansa}

% ================================================================ III
\section{Verificarea cu TeX}

\lettrine{A}{lgoritmul} lui Knuth și Plass e ușor de scris greșit, iar o
greșeală nu se vede: paginile tot arată bine. Ca motorul dispozitivului să poată
fi numit „Knuth–Plass”, el a fost comparat cu TeX, care e implementarea de
referință, pe tot corpusul de cincisprezece cărți.

\subsection{Metoda}
Comparația are sens numai dacă TeX primește exact paragraful pe care l-a avut
dispozitivul. Unealta \fisier{texref.py} reface fiecare paragraf din fișierul
IBV și din fonturi (cutii, spații elastice, cratime moi, ruperi după liniuță)
fără să ia nimic din codul C, în afară de rândurile finale. Lățimile calculate
așa se potrivesc cu ale dispozitivului la 201.015 rânduri din 201.015. În
aceeași unealtă e un model al algoritmului din TeX \cite[§\,813–890]{texprogram},
cu abaterile dispozitivului ca opțiuni. Fără abateri, modelul dă exact rupturile
și penalizarea totală raportate de LuaTeX, la 15.578 de paragrafe din 15.578.
\sursa{tools/texref/ README.md}

\subsection{Prima rulare}
Prima comparație a găsit cinci abateri, niciuna scrisă undeva, și a măsurat-o
pe fiecare singură, peste modelul lui TeX (tabelul~\ref{tab:abateri}). Cea mai
mare era o greșeală: motorul ținea o singură stare pe fiecare punct de rupere,
în loc de câte una pe clasă de potrivire, așa că penalizarea pentru un salt de
la un rând strâns la unul larg era calculată față de un rând care nu fusese
neapărat ales.

\begin{table}[t]
\centering\small
\caption{Abaterile motorului v3 de la TeX: câte paragrafe din prima trecere
schimbă fiecare, pornită singură.}
\label{tab:abateri}
\begin{tabularx}{\textwidth}{@{}>{\itshape}l X >{\lin}r@{}}
\toprule
abaterea & ce face & paragrafe \\
\midrule
one\_node & o singură stare pe punct de rupere, nu câte una pe clasă & 450 \\
fill\_natural & ultimul rând nu se îngustează & 229 \\
exact\_bad & urâțenia exactă, fără plafonul de 10.000 & 74 \\
no\_first\_adj & primul rând nu e comparat cu o clasă dinainte & 58 \\
shrink\_total & îngustarea, rotunjită o dată pe rând & 3 \\
\midrule
toate cinci & & 751 \\
\bottomrule
\end{tabularx}
\end{table}

\subsection{După corectură}
Motorul v4 a păstrat două abateri, cu intenție. A doua trecere a dispozitivului
acceptă orice rând care încape și trebuie să deosebească un rând foarte larg de
unul și mai larg, iar plafonul lui TeX nu permite asta; de aici
\textit{exact\_bad}. A doua, \textit{shrink\_total}, ar cere încă o sumă ținută
pe fiecare punct de rupere pentru câteva unități de 1/64 de pixel. Înainte de a
fi scris, modelul a prezis trecerile motorului nou (12.128 de paragrafe în prima,
3.453 în a doua, 2 fără niciun punct de rupere posibil), iar motorul le-a dat
exact.

\begin{table}[t]
\centering\small
\caption{Motorul dispozitivului față de LuaTeX, pe tot corpusul. Prima trecere
se compară cu TeX la toleranța 400, a doua la 10.000.}
\label{tab:texref}
\lin
\begin{tabular}{@{}lrr@{}}
\toprule
& v3, înainte & v4, după \\
\midrule
prima trecere: paragrafe identice & 93,79 \% & 99,36 \% \\
prima trecere: rânduri identice & 97,60 \% & 99,65 \% \\
a doua trecere: paragrafe identice & 86,84 \% & 96,06 \% \\
toate: paragrafe identice & 92,24 \% & 98,63 \% \\
toate: rânduri identice & 97,59 \% & 99,39 \% \\
\bottomrule
\end{tabular}
\end{table}

Cu metricile din §\,II, motorul corectat are spațiul mediu mai mic sau egal în
toate cele cincisprezece cărți, iar corpusul ocupă cu 20 de pagini mai puțin
(tabelul~\ref{tab:texref}).

% ================================================================ IV
\section{Încăperea cărților}

\lettrine{F}{ormatul IBV} ține textul în blocuri comprimate cu DEFLATE
\cite{rfc1951}, de cel mult 32 KB fiecare. Decompresorul nu ocupă loc în
program: e cel din memoria ROM a lui ESP32-S3. Pentru un salt la o poziție
oarecare, dispozitivul citește și decomprimă un singur bloc.

\subsection{Mărimea}
Textul comprimat are 27–46~\% din textul brut. Toate cărțile ies mai mici decât
sursa, inclusiv EPUB-urile, care sunt deja arhive; cea mai lungă, \textit{Moby
Dick}, are 600 KB. Cratimele moi din §\,II costă: fiecare e un caracter de doi
octeți, iar după compresie adaugă până la 10,2~\% la fișier, la \textit{Țara
mea}. E prețul pentru o despărțire făcută o singură dată, pe calculator, în loc
de la fiecare paginare.
\sursa{measurements/E5/ REZULTATE.md}

\subsection{Blocul}
Un bloc mai mare se comprimă mai bine, dar un salt îl citește întreg, iar
memoria pentru decomprimare crește cu el (tabelul~\ref{tab:bloc}). Câștigul de
5,7~\% al blocurilor de 32 KB înseamnă vreo 100 KB pe toate cele opt cărți în
proză, nimic pe un card de câțiva gigaocteți; de aceea rămâne 16 KB. Cât costă
de fapt un salt pe dispozitiv nu e măsurat.

\begin{table}[t]
\centering\small
\caption{Cărțile în proză, pe mărimea țintă a blocului. Mărimea e a tuturor
celor opt fișiere, față de blocurile de 16 KB.}
\label{tab:bloc}
\lin
\ceva{date/blocuri.dat}{%
\pgfplotstabletypeset[
  columns={kb,relativ,citit_kb},
  columns/kb/.style={column name={bloc (KB)}, column type={@{}r}},
  columns/relativ/.style={column name={mărimea}, fixed, precision=1, showpos, zerofill,
    use comma, postproc cell content/.append style={/pgfplots/table/@cell content/.add={}{\,\%}}},
  columns/citit_kb/.style={column name={citit la un salt, mediana (KB)}, fixed, precision=1,
    zerofill, use comma, column type={r@{}}},
  every head row/.style={before row=\toprule, after row=\midrule},
  every last row/.style={after row=\bottomrule}]{date/blocuri.dat}}{\lipseste{0.6\textwidth}{20mm}}
\sursa{E5/date/ blocuri.csv}
\end{table}

\subsection{Deschiderea}
La nouă cărți din douăsprezece, tot ce trebuie pentru deschidere (antetul,
directorul, metadatele, cuprinsul și indexul blocurilor) e în primii 4 KB ai
fișierului. Celelalte trei au cuprinsuri lungi: \textit{Moby Dick}, Eminescu și
Whitman. Conversia unei cărți durează 0,23–1,67 secunde pe un laptop.

% ================================================================ V
\section{Carcasa și curentul}

\lettrine{P}{laca proprie} are exact conturul modulului de afișaj de 3,7 țoli,
105,79 × 53,80 mm, și stă sub el, în aceeași stivă de șuruburi. Are patru
straturi și trece verificarea regulilor electrice și de desen fără nicio eroare.
Bateria stă dedesubt, într-un buzunar al capacului. Planșa~III e o secțiune
prin mijlocul carcasei, desenată din aceleași cote din care e generată carcasa
pentru tipărire.
\sursa{hardware/pcb/ README.md, enclosure/ carcasa.scad}

\begin{plansa}[t]
\makebox[\textwidth][l]{%
\ceva{stiva.mps}{%
\begin{tikzpicture}
  \node[anchor=south west, inner sep=0] at (0,0) {\includegraphics[width=\latplin]{stiva.mps}};
  \begin{scope}[x=\latplin/122.79, y=\latplin/122.79, shift={(4.5,13)}]
    \lit{a}{78,-2.9}{75,0.8}
    \lit{b}{-7.5,16}{-1.4,16}
    \node[fill=white, inner sep=0.8pt, font=\small\itshape] at (44,4.6) {c};
    \lit{d}{30,13}{30,9.4}
    \lit{e}{76,15}{86,12.4}
    \lit{f}{50,18.2}{50,21.4}
    \lit{g}{62,29}{62,23.4}
    \lit{h}{30,29}{30,24.2}
    \lit{i}{13.5,4.6}{8.4,4.6}
    \lit{k}{13.5,16}{8.0,16}
    \lit{l}{104,29}{99.2,24.4}
    \node[fill=white, inner sep=1pt, font=\footnotesize\lin] at (52.9,-5.5) {110,8};
    \node[fill=white, inner sep=1pt, font=\footnotesize\lin, rotate=90] at (114.29,12.9) {25,8};
    \node[anchor=west, font=\footnotesize\lin] at (8,-10.3) {10 mm};
  \end{scope}
\end{tikzpicture}}{\lipseste{\latplin}{60mm}}}
\caption[Secțiune prin carcasă]{Secțiune prin carcasă, în planul de mijloc al
plăcii; cotele în milimetri. \textit{a}, capacul; \textit{b}, rama, cu pereții;
\textit{c}, bateria, 50 × 34 × 6 mm (presupusă); \textit{d}, placa proprie;
\textit{e}, modulul ESP32-S3-WROOM-1, cu antena la margine; \textit{f}, placa
modulului de afișaj; \textit{g}, sticla panoului; \textit{h}, zona activă,
81,40 mm; \textit{i}, stâlpii capacului; \textit{k}, distanțierele dintre plăci;
\textit{l}, șuruburile M2 × 25 cu cap înecat și piulițele lor, în spatele
planului, deci întrerupte. Piesele din plan sunt hașurate; cele din spatele lui,
umbrite.}
\end{plansa}

\subsection{Somnul}
Cât timp o pagină e citită, cititorul doarme. Comutatoarele de sarcină taie
alimentarea cardului și a panoului, iar ce rămâne e în tabelul~\ref{tab:somn}.
Împărțind capacitatea bateriei la acest curent,
\begin{equation}
  \frac{1500~\text{mAh}}{17{,}5~\text{µA}} \approx 85\,700~\text{h} \approx 9{,}8~\text{ani},
\end{equation}
ceea ce nu e o durată de funcționare, ci doar arată că somnul nu e problema:
o baterie LiPo se descarcă singură mai repede de atât. Ce golește bateria sunt
paginile întoarse, și acelea se măsoară.

\begin{table}[t]
\centering\small
\caption{Curentul în somn, pe piese: valori tipice din foile de catalog,
nemăsurate.}
\label{tab:somn}
\begin{tabular}{@{}l>{\lin}r@{}}
\toprule
piesa & µA \\
\midrule
ESP32-S3, somn adânc, cu temporizatorul RTC & \textasciitilde 8 \\
BQ24074, din baterie, fără USB & 4,3 \\
MAX17048, în hibernare & 3 \\
divizorul tensiunii bateriei, 2 M\ohm & 2,1 \\
TPS63900, în gol & \textasciitilde 0,1 \\
cardul și panoul, cu comutatoarele oprite & \textasciitilde 0 \\
\midrule
total & \textasciitilde 17,5 \\
\bottomrule
\end{tabular}
\sursa{hardware/pcb/ README.md}
\end{table}

\subsection{Măsurarea}
Pentru energia pe pagină, un INA226 cu un șunt de 0,1~\ohm stă în serie cu
bateria, iar o a doua placă îl citește la fiecare conversie, cam de 850 de ori
pe secundă. Rezoluția curentului este
\begin{equation}
  \Delta I = \frac{2{,}5~\text{µV}}{0{,}1~\text{\ohm}} = 25~\text{µA},
\end{equation}
prea grosieră pentru somn, care se măsoară cu multimetrul, dar suficientă
pentru o reîmprospătare a ecranului. Energia unei pagini este integrala puterii
pe intervalul întoarcerii, prin trapeze:
\begin{equation}
  E = \int_{t_0}^{t_1} V I \,\mathrm{d}t
    \approx \sum_k \frac{(t_{k+1} - t_k)\,(V_k I_k + V_{k+1} I_{k+1})}{2}.
\end{equation}
Intervalele vin din jurnalul de timp al cititorului, care are alt ceas;
unealta le aliniază căutând decalajul care pune toate întoarcerile peste
curentul cel mai mare. Pe date sintetice găsește decalajul exact.
\sursa{tools/ina226/ README.md}

% ================================================================ VI
\section{Ce nu e dovedit}

O lucrare care măsoară trebuie să spună la fel de limpede și ce
n-a măsurat. Iată tot ce presupune textul de până aici fără să fi văzut cu
ochii pe un dispozitiv adevărat:
\begin{itemize}[leftmargin=1.4em]
  \item nimic nu a rulat pe hardware; testele dispozitivului (22, 29 și 41 de
    cazuri, toate trecute) rulează într-un emulator;
  \item placa nu a fost comandată, deci nici pornită; carcasa nu a fost
    tipărită, iar bateria are dimensiuni presupuse;
  \item consumul în somn e o sumă de valori tipice din cataloage;
  \item timpul de paginare e măsurat pe un laptop, nu pe ESP32-S3; la fel
    costul unui salt în carte;
  \item măsurarea energiei pe pagină e scrisă și testată pe date sintetice,
    dar modulul de măsură nu a fost cumpărat.
\end{itemize}
Ce e dovedit: paginile din Planșa~II sunt cele ale codului dispozitivului,
algoritmul lui de rupere e cel al lui TeX cu două abateri numite, iar cifrele
din §§\,II și IV sunt reproductibile din fișierele depozitului.

% ================================================================ bibliografia
\renewcommand\refname{\normalfont\large\itshape Bibliografie}
\begin{thebibliography}{9}\small
\bibitem{knuthplass} \textsc{D. E. Knuth, M. F. Plass}, „Breaking paragraphs
  into lines”, \textit{Software: Practice and Experience} 11(11), 1981,
  p.~1119–1184.
\bibitem{texprogram} \textsc{D. E. Knuth}, \textit{TeX: The Program}, Computers
  \& Typesetting, vol.~B, Addison-Wesley, 1986.
\bibitem{liang} \textsc{F. M. Liang}, \textit{Word Hy-phen-a-tion by
  Com-put-er}, teză de doctorat, Stanford University, 1983.
\bibitem{comiskey} \textsc{B. Comiskey, J. D. Albert, H. Yoshizawa,
  J. Jacobson}, „An electrophoretic ink for all-printed reflective electronic
  displays”, \textit{Nature} 394, 1998, p.~253–255.
\bibitem{rfc1951} \textsc{P. Deutsch}, \textit{DEFLATE Compressed Data Format
  Specification version 1.3}, RFC 1951, 1996.
\bibitem{ina226} \textsc{Texas Instruments}, \textit{INA226: High-Side or
  Low-Side Measurement, Bi-Directional Current and Power Monitor with
  I²C Compatible Interface}, foaie de catalog.
\bibitem{fiziko} \textsc{S. Slyusarev}, \textit{fiziko}, bibliotecă MetaPost
  pentru ilustrații în stilul manualelor vechi, GPL-3.0.
\bibitem{gutenberg} \textit{Project Gutenberg}, sursa cărților din corpus.
\end{thebibliography}

% ================================================================ colofon
\vfill
\begin{center}
\ornament\\[3mm]
\begin{minipage}{0.7\textwidth}\centering\footnotesize\itshape
Cules în Libron, de Nico Verbruggen (licența OFL), fontul cu care cititorul
însuși arată textul latin. Așezat cu LuaLaTeX; paragrafele sunt rupte cu
parametrii dispozitivului și despărțite după tiparele convertorului. Planșele
sunt gravate în MetaPost, cu fiziko de Sergey Slyusarev. Figurile și tabelele
din §§\,II și IV sunt citite la compilare din datele măsurătorilor.\\[2mm]
{\upshape\lin Construit din \versiune.}
\end{minipage}
\end{center}
\end{document}
